% Generado por scripts/generate_resumen_ejecutivo.py. No editar a mano.

\begingroup
\small
\arrayrulecolor{Azul-UCR!80}%
\renewcommand{\arraystretch}{1.2}%
\rowcolors{2}{white}{Celeste-UCR!12}%
\begin{longtable}{|>{\raggedright\arraybackslash}p{3.3cm}|>{\raggedright\arraybackslash}p{5.5cm}|>{\raggedright\arraybackslash}p{6.9cm}|}
\caption{Descripción de los cursos propios de la carrera}\label{res:descripciones}\\
\hline
\rowcolor{Azul-UCR}
\textcolor{white}{\textbf{Curso (créditos)}} & \textcolor{white}{\textbf{Objetivo general}} & \textcolor{white}{\textbf{Temática resumida}} \\ \hline
\endfirsthead
\hline
\rowcolor{Azul-UCR}
\textcolor{white}{\textbf{Curso (créditos)}} & \textcolor{white}{\textbf{Objetivo general}} & \textcolor{white}{\textbf{Temática resumida}} \\ \hline
\endhead
\hline
\endfoot
\textbf{MA-0151}\newline Fundamentos de Álgebra, Trigonometría y Geometría Analítica (3) & Aplicar las nociones, intuiciones y conceptos básicos del álgebra, trigonometría y geometría analítica para la resolución de problemas. & Números reales y sus propiedades básicas; Trigonometría; Números complejos; Geometría vectorial y analítica básica; Proyectos de descubrimiento. \\ \hline
\textbf{MA-0152}\newline Matemática Exploratoria (3) & Iniciar al estudiantado en el quehacer matemático mediante la resolución de problemas de teoría de números, combinatoria y geometría, a través de la exploración, la conjetura y la justificación de resultados. & Introducción a los modos de trabajo y desarrollo de la matemática; Estrategias generales de resolución de problemas; Resolución de problemas concretos; Introducción a la formalización de cuantificadores. \\ \hline
\textbf{MA-0251}\newline Cálculo en una Variable (4) & Aplicar las nociones, intuiciones y conceptos básicos del cálculo diferencial e integral en una variable para la resolución de problemas. & Introducción a las funciones; Límites y continuidad; Derivación; Integración de Riemann. \\ \hline
\textbf{MA-0252}\newline Introducción a las Demostraciones (3) & Desarrollar en el estudiantado las habilidades de comprensión y aplicación de estrategias y métodos de demostración en matemáticas. & Lógica básica y estrategias de demostración; Teoría de conjuntos; Relaciones y funciones; Inducción matemática y recurrencia; Teoría de números. \\ \hline
\textbf{MA-0341}\newline Matemática Computacional (4) & Implementar el pensamiento algorítmico en la resolución de problemas matemáticos, mediante la ejecución de programas en lenguajes de programación de alto nivel. & Introducción a lenguajes de programación; Introducción a algoritmos; Visualizaciones; Algoritmos básicos en teoría de números; Algoritmos básicos en cálculo diferencial e integral; Algoritmos básicos en álgebra lineal y aplicaciones. \\ \hline
\textbf{MA-0351}\newline Cálculo en Varias Variables (4) & Aplicar las nociones, intuiciones y conceptos básicos del cálculo diferencial e integral en varias variables para la resolución de problemas. & Sucesiones y series; Derivación de funciones de varias variables; Integración en varias variables operacional. \\ \hline
\textbf{MA-0361}\newline Álgebra Lineal I (4) & Aplicar las nociones, intuiciones y conceptos básicos del álgebra lineal para la resolución de problemas. & Sistemas de ecuaciones lineales; Vectores y matrices; Determinantes; Espacios vectoriales; Transformaciones lineales. \\ \hline
\textbf{MA-0451}\newline Principios de Análisis I (4) & Aplicar los conceptos básicos de completitud, convergencia y diferenciación para la demostración de propiedades del análisis en una variable real. & Completitud y propiedades de los números reales; Sucesiones y series numéricas; Límites y continuidad; Diferenciación; Sucesiones y series de funciones; Construcción de las funciones trascendentes elementales. \\ \hline
\textbf{MA-0461}\newline Álgebra Lineal II (4) & Aplicar las nociones, intuiciones y conceptos del álgebra lineal para la resolución de problemas. & Polinomios; Dualidad; Valores y vectores propios; Espacios con producto interno; Triangularización, diagonalización y formas canónicas; Formas bilineales y multilineales. \\ \hline
\textbf{MA-0471}\newline Introducción a la Geometría Diferencial (3) & Analizar de forma rigurosa e intuitiva la geometría diferencial de curvas y superficies, mediante cálculos y demostraciones sobre sus propiedades locales y globales, tanto extrínsecas como intrínsecas. & Curvas; Superficies; Aplicación de Gauss; Geometría intrínseca de superficies; Geometría proyectiva. \\ \hline
\textbf{MA-0496}\newline Teoría de Números (3) & Aplicar los conceptos básicos de la teoría de números elemental para la resolución de problemas. & Repaso de divisibilidad en los enteros; Números primos; Congruencias; Funciones aritméticas; Raíces primitivas; Congruencias cuadráticas y la Ley de Reciprocidad cuadrática; Enteros en cuerpos de números cuadráticos. \\ \hline
\textbf{MA-0501}\newline Análisis Numérico I (4) & Aplicar la teoría de análisis numérico y sus herramientas básicas en la resolución de problemas y modelos matemáticos, mediante el estudio teórico de algoritmos y la implementación de algoritmos en un lenguaje de alto nivel. & Error numérico; Aproximación de funciones; Ecuaciones no lineales; Integración; Álgebra lineal numérica; Ecuaciones diferenciales ordinarias; Tema opcional. \\ \hline
\textbf{MA-0556}\newline Ecuaciones Diferenciales (4) & Establecer criterios para determinar la existencia, unicidad y estabilidad de soluciones de ecuaciones diferenciales ordinarias. & Introducción y ecuaciones diferenciales ordinarias de primer orden; Sistemas de ecuaciones diferenciales lineales; La transformada de Laplace; Teoría de existencia y unicidad; Continuación de soluciones; Dependencia de valores iniciales y parámetros; Tema a elegir por la persona docente. \\ \hline
\textbf{MA-0555}\newline Principios de Análisis II (4) & Aplicar los conceptos básicos de convergencia, diferenciación e integración para la demostración de propiedades del análisis en una y varias variables reales. & Integración de Riemann en $\mathbb{R}$; Topología métrica de $\mathbb{R}^n$; Funciones de varias variables; Diferenciación en varias variables; Cálculo y análisis vectorial. \\ \hline
\textbf{MA-0635}\newline Topología General (4) & Comprender los conceptos fundamentales de los espacios métricos y de los espacios topológicos, así como las funciones continuas entre ellos, y desarrollar herramientas, tanto conjuntistas como algebraicas, para analizar y distinguir dichos espacios. & Repaso de teoría de conjuntos; Espacios métricos; Espacios topológicos; Axiomas de separación; Homotopía y el grupo fundamental. \\ \hline
\textbf{MA-0571}\newline Álgebra Abstracta I (4) & Introducir a la persona estudiante a los conceptos y métodos propios del álgebra abstracta a través del estudio de la teoría de anillos y la teoría de módulos. & Definiciones, ejemplos y propiedades básicas de anillos; Morfismos, ideales y cocientes de anillos; Teoría de divisibilidad y factorización en anillos conmutativos; Anillos de polinomios y factorización; Definiciones, ejemplos y propiedades básicas de módulos; Grupos. \\ \hline
\textbf{MA-0661}\newline Álgebra Abstracta II (5) & Analizar y aplicar los conceptos fundamentales de la teoría de grupos, cuerpos y extensiones de cuerpos, para comprender la estructura de las extensiones algebraicas y su relación con los grupos de automorfismos, mediante el uso de resultados clásicos de álgebra. & Teoremas de estructura de grupos; Cuerpos; Automorfismos; Extensiones normales; Extensiones separables e inseparables; Teorema Fundamental de la Teoría de Galois; Cuerpos finitos. \\ \hline
\textbf{MA-0647}\newline Modelación Matemática (4) & Modelar problemas determinísticos a través de ecuaciones diferenciales y sistemas dinámicos para situaciones de ciencias, ingeniería o contextos interdisciplinarios, evaluando la validez del modelo y comunicando resultados con claridad. & El ciclo de modelación matemática; Modelación continua con ecuaciones diferenciales y sistemas dinámicos; Modelación discreta y simulación computacional; Validación, incertidumbre e interpretación con datos; Temas complementarios a elección de la persona docente; Proyecto final integrador (transversal, semanas 11--15). \\ \hline
\textbf{MA-0730}\newline Análisis Real I (5) & Aplicar los conceptos y resultados básicos del análisis real para la resolución de problemas. & La integral de Riemann-Stieltjes; Medida de Lebesgue en $\R^n$; Funciones Lebesgue medibles; La integral de Lebesgue; Espacios $L^p$. \\ \hline
\textbf{MA-0735}\newline Análisis Complejo (4) & Aplicar los conceptos y resultados básicos del análisis complejo para la resolución de problemas teóricos y prácticos. & Números complejos y funciones en el plano complejo; Diferenciación en el plano complejo; Integración compleja, el Teorema de Cauchy y sus aplicaciones; Singularidades, funciones meromorfas y el logaritmo complejo; Funciones enteras; Tópicos adicionales (cubrir al menos uno a criterio de la persona docente). \\ \hline
\textbf{MA-0880}\newline Comunicación en las Ciencias (2) & Que la persona estudiante sea capaz de comunicarse en forma oral y escrita para transmitir de forma efectiva ideas, resultados y argumentos matemáticos con claridad, rigor y pertinencia en contextos académicos y profesionales. & Comunicación científica y ética profesional; Escritura matemática y científica; Presentación oral y apoyos visuales; Búsqueda bibliográfica y gestión de fuentes; Comunicación en docencia y divulgación; Planificación de carrera y comunicación profesional; Presentaciones integradoras del estudiantado. \\ \hline
\textbf{MA-0881}\newline Práctica Profesional en Matemática Pura (6) & Que la persona estudiante integre y aplique conocimientos de matemática pura en un entorno académico o profesional, desarrollando autonomía, responsabilidad y capacidad de comunicar resultados. & Los contenidos no se fijan de manera permanente en el plan de estudios: dependen de la modalidad y del plan de trabajo aprobado para cada persona estudiante. \\ \hline
\textbf{MA-0883}\newline Práctica Profesional en Matemática Aplicada (6) & Que la persona estudiante integre y aplique conocimientos de matemática aplicada en un entorno académico o profesional, desarrollando autonomía, responsabilidad y capacidad de comunicar resultados. & Los contenidos no se fijan de manera permanente en el plan de estudios: dependen de la modalidad y del plan de trabajo aprobado para cada persona estudiante. \\ \hline
\textbf{MA-0982}\newline Seminario de Estudios Dirigidos en Matemática Pura I (7) & Que la persona estudiante desarrolle, de manera autónoma, actividades asociadas a la investigación en matemática pura. & Los contenidos no se fijan de manera permanente en el plan de estudios: dependen del tema y de los objetivos acordados para cada persona estudiante. \\ \hline
\textbf{MA-1081}\newline Seminario de Estudios Dirigidos en Matemática Pura II (7) & Que la persona estudiante desarrolle, de manera autónoma, actividades asociadas a la investigación en matemática pura. & Los contenidos no se fijan de manera permanente en el plan de estudios: dependen del tema y de los objetivos acordados para cada persona estudiante. \\ \hline
\textbf{MA-0984}\newline Seminario de Estudios Dirigidos en Matemática Aplicada I (7) & Que la persona estudiante desarrolle, de manera autónoma, actividades asociadas a la investigación en matemática aplicada. & Los contenidos no se fijan de manera permanente en el plan de estudios: dependen del tema y de los objetivos acordados para cada persona estudiante. \\ \hline
\textbf{MA-1082}\newline Seminario de Estudios Dirigidos en Matemática Aplicada II (7) & Que la persona estudiante desarrolle, de manera autónoma, actividades asociadas a la investigación en matemática aplicada. & Los contenidos no se fijan de manera permanente en el plan de estudios: dependen del tema y de los objetivos acordados para cada persona estudiante. \\ \hline
\textbf{MA-0526}\newline Combinatoria (4) & Desarrollar en la persona estudiante competencias para abordar problemas combinatorios mediante técnicas de conteo, teoría de grafos, métodos probabilísticos y combinatoria extremal. & Herramientas de conteo; Introducción a la teoría de grafos; Método probabilístico; Funciones generatrices; Combinatoria extremal. \\ \hline
\textbf{MA-0528}\newline Sistemas Dinámicos (4) & Que la persona estudiante adquiera las herramientas fundamentales para analizar sistemas dinámicos desde los puntos de vista topológico, diferenciable y ergódico, reconocer estructuras recurrentes, construir modelos simbólicos, estudiar medidas invariantes e interpretar el comportamiento estadístico de las órbitas mediante promedios temporales. & Sistemas dinámicos en tiempo discreto y continuo; Conjuntos límite, recurrencia y conjugaciones; Difeomorfismos del intervalo y linealización local; Codificación simbólica; Medidas invariantes y teoría ergódica; Mapas por ramas, mapa de Gauss y distorsión acotada; La familia cuadrática. \\ \hline
\textbf{MA-0838}\newline Geometría Diferencial (4) & Comprender y aplicar las herramientas fundamentales del cálculo diferencial e integral y de la geometría riemanniana en variedades, demostrando resultados centrales. & Variedades diferenciales y topología diferencial; Formas diferenciales; Integración en variedades; Geometría riemanniana. \\ \hline
\textbf{MA-0837}\newline Geometría Algebraica (4) & Introducir a la persona estudiante a los conceptos, métodos y resultados fundamentales de la Geometría Algebraica clásica, con énfasis en variedades y curvas algebraicas, sus singularidades, propiedades proyectivas y birracionales, y sus invariantes geométricos y aritméticos. \ & Conjuntos algebraicos afines y el Teorema de los Ceros de Hilbert; Variedades afines y propiedades locales; Propiedades locales de curvas planas; Geometría proyectiva; Curvas proyectivas planas e intersecciones; Morfismos y geometría birracional; Resolución de singularidades de curvas; Divisores, género y Teorema de Riemann--Roch; Curvas sobre cuerpos finitos, funciones zeta y conjeturas de Weil. \\ \hline
\textbf{MA-0734}\newline Optimización (4) & Desarrollar los fundamentos matemáticos de la optimización para establecer la existencia y caracterizar las soluciones de problemas con y sin restricciones. & Fundamentos y optimización sin restricciones; Programación lineal; Optimización con restricciones de igualdad y desigualdad; Proyectos de descubrimiento. \\ \hline
\textbf{MA-0756}\newline Ecuaciones Diferenciales Parciales (5) & Desarrollar competencias teóricas para formular, analizar y resolver problemas de ecuaciones diferenciales parciales mediante métodos clásicos y formulaciones débiles. & Ecuaciones clásicas fundamentales; Espacios de Sobolev; Ecuaciones elípticas de segundo orden; Ecuaciones de evolución lineales. \\ \hline
\textbf{MA-0781}\newline Lógica (5) & Desarrollar los conocimientos y herramientas fundamentales de la lógica matemática para analizar y formalizar razonamientos matemáticos y demostrar los principales resultados del cálculo de predicados de primer orden. & Ordinales y cardinales; Lógica de primer orden; Teorema de completitud y teorema de compacidad; Nociones básicas de teorías y estructuras. \\ \hline
\textbf{MA-0783}\newline Integración (4) & Desarrollar competencias en teoría abstracta de la medida e integración y en sus aplicaciones clásicas dentro del análisis matemático. & Repaso de integración de Lebesgue; Curvas rectificables e isoperimetría; Teoría abstracta de la medida e integración; Teoremas ergódicos; Medidas de Hausdorff; Medidas de Haar; Transporte óptimo. \\ \hline
\textbf{MA-0788}\newline Análisis Real II (4) & Aplicar los conceptos y resultados intermedios del Análisis Real para la resolución de problemas. & Diferenciación; Continuidad absoluta y el Teorema Fundamental del Cálculo; Espacios de Banach; Espacios de Hilbert; Introducción al Análisis de Fourier. \\ \hline
\textbf{MA-0835}\newline Topología Algebraica (4) & Que la persona estudiante domine los conceptos y técnicas elementales de la topología algebraica y pueda aplicarlos al análisis de espacios y aplicaciones continuas. & Homotopía (noción básica); Grupo fundamental; Homología; Cohomología de haces (Sheaf Cohomology); Homotopía superior. \\ \hline
\textbf{MA-0834}\newline Análisis Funcional (4) & Que la persona estudiante sea capaz de resolver problemas y desarrollar teoría a partir de los fundamentos del análisis funcional en espacios de Banach y de Hilbert. & Espacios de Banach y espacios $L^p$; Espacios de Hilbert; Operadores lineales acotados; Teoremas fundamentales; Dualidad y topologías débiles; Operadores compactos y teorema espectral. \\ \hline
\textbf{MA-0863}\newline Teoría Analítica de Números (4) & Introducir a la persona estudiante a las principales ideas, técnicas y resultados de la Teoría Analítica de Números, enfatizando el uso de herramientas del análisis real y complejo para estudiar funciones aritméticas y la distribución de los números primos. \ & Funciones aritméticas y convolución de Dirichlet; Funciones sumatorias y órdenes promedio; Series de Dirichlet y herramientas de análisis complejo; La función zeta de Riemann; Ceros de la función zeta y distribución de los números primos; Caracteres y funciones $L$ de Dirichlet; Primos en progresiones aritméticas; Problemas abiertos y desarrollos posteriores. \\ \hline
\textbf{MA-0862}\newline Análisis Armónico (4) & Que la persona estudiante sea capaz de resolver problemas y desarrollar teoría a partir de los fundamentos del análisis armónico de variable real en $\R^n$. & Espacios $L^p$; La transformada de Fourier; Funciones maximales; Las transformadas de Hilbert y de Riesz; Integrales singulares de tipo convolución; Teoría de Littlewood--Paley. \\ \hline
\textbf{MA-0861}\newline Teoría de Modelos (4) & Comprender los principios básicos de la teoría de modelos y de sus aplicaciones a las estructuras algebraicas. & Fundamentos de lógica de primer orden; Principios básicos de teoría de modelos; Eliminación de cuantificadores; Aplicaciones al álgebra y la geometría algebraica; Tipos y compacidad. \\ \hline
\textbf{MA-0841}\newline Probabilidad (4) & Desarrollar los fundamentos de la teoría de la probabilidad en el marco de la teoría de la medida para modelar incertidumbre, demostrar resultados centrales y aplicarlos en problemas matemáticos e interdisciplinarios. & Probabilidad elemental; Teoría de la medida e integración; Teoría de probabilidad; Independencia; Convergencia de variables aleatorias; Probabilidad condicional. \\ \hline
\textbf{MA-0849}\newline Teoría de Módulos (5) & Comprender los fundamentos de la teoría de módulos y sus principales aplicaciones algebraicas. & Módulos; Construcciones básicas; Módulos noetherianos y artinianos; Módulos sobre dominios de ideales principales; Módulos proyectivos, inyectivos y planos; Producto tensorial; Introducción al álgebra homológica (Optativo). \\ \hline
\textbf{MA-0858}\newline Teoría de Conjuntos (4) & Introducir al estudiantado en la teoría axiomática de conjuntos de Zermelo--Fraenkel, desarrollando las herramientas fundamentales como cardinales, ordinales y aritmética transfinita. & Motivación y antecedentes históricos; El sistema axiomático de Zermelo--Fraenkel; Relaciones y funciones; Órdenes; Los números naturales; Números cardinales; Números ordinales; El axioma de elección y la hipótesis del continuo. \\ \hline
\textbf{MA-0876}\newline Teoría Algebraica de Números (4) & Introducir a la persona estudiante a los conceptos, métodos y ejemplos fundamentales de la Teoría de Números Algebraicos, enfatizando la aritmética de los anillos de enteros de cuerpos de números, la factorización de ideales, el grupo de clases, el grupo de unidades, la factorización de primos en extensiones y el uso de herramientas computacionales para explorar ejemplos concretos. & Cuerpos de números y enteros algebraicos; Ideales y factorización; Factorización de ideales primos y ramificación; El grupo de clases y el teorema de Minkowski; Unidades y teorema de Dirichlet; Conteo de ideales y función zeta de Dedekind; Teoría de descomposición e inercia en extensiones de Galois. \\ \hline
\textbf{MA-0889}\newline Álgebra Conmutativa (4) & Introducir a la persona estudiante a los conceptos, métodos y resultados fundamentales del Álgebra Conmutativa, enfatizando las propiedades de anillos y módulos noetherianos, la localización, la dependencia integral, la teoría de dimensión, los métodos computacionales y la estructura local de los anillos conmutativos. & Módulos y herramientas algebraicas fundamentales; El diccionario entre álgebra y geometría; Noetherianidad y descomposición primaria; Localización y anillos locales; Dependencia integral y normalización; Dimensión de Krull; Bases de Gröbner y métodos computacionales; Anillos locales regulares y anillos de dimensión uno. \\ \hline
\textbf{MA-0920}\newline Ecuaciones en Derivadas Parciales Numéricas (4) & Aplicar la teoría del análisis numérico y sus herramientas básicas para aproximar soluciones numéricas de ecuaciones diferenciales parciales. & Introducción a las ecuaciones diferenciales parciales; Método de diferencias finitas; Método de elemento finito; Resolución numérica de sistemas lineales; Resolución numérica de la ecuación de calor y de onda; Temas adicionales. \\ \hline
\textbf{MA-0477}\newline Ecuaciones Diferenciales Parciales Aplicadas (4) & Aplicar métodos clásicos del análisis de EDP para resolver y analizar ecuaciones lineales que surgen en modelos matemáticos. & EDP de primer orden; Ecuaciones de Laplace y Poisson; Ecuación del calor; Ecuación de onda; Introducción al cálculo de variaciones. \\ \hline
\textbf{MA-0759}\newline Tópicos de Matemática Financiera (4) & Profundizar en un tema de la matemática financiera, definido para cada oferta del curso, mediante el estudio guiado de resultados, problemas y, cuando corresponda, aplicaciones. & Los contenidos no se fijan de manera permanente en el plan de estudios: son abiertos por oferta, siempre que pertenezcan al ámbito de Matemática financiera. Al inicio del ciclo lectivo, la persona docente debe comunicar por escrito al estudiantado. \\ \hline
\textbf{MA-0790}\newline Tópicos de Análisis (4) & Profundizar en un tema de análisis matemático, definido para cada oferta del curso, mediante el estudio guiado de resultados, problemas y, cuando corresponda, aplicaciones. & Los contenidos no se fijan de manera permanente en el plan de estudios: son abiertos por oferta, siempre que pertenezcan al ámbito del Análisis matemático. \\ \hline
\textbf{MA-0791}\newline Tópicos de Probabilidad (4) & Profundizar en un tema de la probabilidad, definido para cada oferta del curso, mediante el estudio guiado de resultados, problemas y, cuando corresponda, aplicaciones. & Los contenidos no se fijan de manera permanente en el plan de estudios: son abiertos por oferta, siempre que pertenezcan al ámbito de Probabilidad. Al inicio del ciclo lectivo, la persona docente debe comunicar por escrito al estudiantado. \\ \hline
\textbf{MA-0792}\newline Tópicos de Ecuaciones Diferenciales (4) & Profundizar en un tema de las ecuaciones diferenciales, definido para cada oferta del curso, mediante el estudio guiado de resultados, problemas y, cuando corresponda, aplicaciones. & Los contenidos no se fijan de manera permanente en el plan de estudios: son abiertos por oferta, siempre que pertenezcan al ámbito de Ecuaciones diferenciales. Al inicio del ciclo lectivo, la persona docente debe comunicar por escrito al estudiantado. \\ \hline
\textbf{MA-0793}\newline Tópicos de Geometría (4) & Profundizar en un tema de la geometría, definido para cada oferta del curso, mediante el estudio guiado de resultados, problemas y, cuando corresponda, aplicaciones. & Los contenidos no se fijan de manera permanente en el plan de estudios: son abiertos por oferta, siempre que pertenezcan al ámbito de Geometría. Al inicio del ciclo lectivo, la persona docente debe comunicar por escrito al estudiantado. \\ \hline
\textbf{MA-0794}\newline Tópicos de Modelación (4) & Profundizar en un tema de la modelación matemática, definido para cada oferta del curso, mediante el estudio guiado de resultados, problemas y, cuando corresponda, aplicaciones. & Los contenidos no se fijan de manera permanente en el plan de estudios: son abiertos por oferta, siempre que pertenezcan al ámbito de Modelación matemática. Al inicio del ciclo lectivo, la persona docente debe comunicar por escrito al estudiantado. \\ \hline
\textbf{MA-0795}\newline Tópicos de Computación Científica (4) & Profundizar en un tema de la computación científica, definido para cada oferta del curso, mediante el estudio guiado de resultados, problemas y, cuando corresponda, aplicaciones. & Los contenidos no se fijan de manera permanente en el plan de estudios: son abiertos por oferta, siempre que pertenezcan al ámbito de Computación científica. Al inicio del ciclo lectivo, la persona docente debe comunicar por escrito al estudiantado. \\ \hline
\textbf{MA-0796}\newline Tópicos de Álgebra (4) & Profundizar en un tema del álgebra, definido para cada oferta del curso, mediante el estudio guiado de resultados, problemas y, cuando corresponda, aplicaciones. & Los contenidos no se fijan de manera permanente en el plan de estudios: son abiertos por oferta, siempre que pertenezcan al ámbito de Álgebra. Al inicio del ciclo lectivo, la persona docente debe comunicar por escrito al estudiantado. \\ \hline
\textbf{MA-0797}\newline Tópicos de Matemática Discreta (4) & Profundizar en un tema de la matemática discreta, definido para cada oferta del curso, mediante el estudio guiado de resultados, problemas y, cuando corresponda, aplicaciones. & Los contenidos no se fijan de manera permanente en el plan de estudios: son abiertos por oferta, siempre que pertenezcan al ámbito de Matemática discreta. Al inicio del ciclo lectivo, la persona docente debe comunicar por escrito al estudiantado. \\ \hline
\textbf{MA-0798}\newline Tópicos de Análisis de Datos (4) & Profundizar en un tema del análisis de datos, definido para cada oferta del curso, mediante el estudio guiado de resultados, problemas y, cuando corresponda, aplicaciones. & Los contenidos no se fijan de manera permanente en el plan de estudios: son abiertos por oferta, siempre que pertenezcan al ámbito de Análisis de datos. Al inicio del ciclo lectivo, la persona docente debe comunicar por escrito al estudiantado. \\ \hline
\textbf{MA-0836}\newline Tópicos de Aprendizaje Automático (4) & Profundizar en un tema del aprendizaje automático, definido para cada oferta del curso, mediante el estudio guiado de resultados, problemas y, cuando corresponda, aplicaciones. & Los contenidos no se fijan de manera permanente en el plan de estudios: son abiertos por oferta, siempre que pertenezcan al ámbito de Aprendizaje automático. Al inicio del ciclo lectivo, la persona docente debe comunicar por escrito al estudiantado. \\ \hline
\textbf{MA-0830}\newline Tópicos de Teoría de Números (4) & Profundizar en un tema de la teoría de números, definido para cada oferta del curso, mediante el estudio guiado de resultados, problemas y, cuando corresponda, aplicaciones. & Los contenidos no se fijan de manera permanente en el plan de estudios: son abiertos por oferta, siempre que pertenezcan al ámbito de Teoría de números. Al inicio del ciclo lectivo, la persona docente debe comunicar por escrito al estudiantado. \\ \hline
\textbf{MA-0831}\newline Tópicos de Lógica (4) & Profundizar en un tema de la lógica, definido para cada oferta del curso, mediante el estudio guiado de resultados, problemas y, cuando corresponda, aplicaciones. & Los contenidos no se fijan de manera permanente en el plan de estudios: son abiertos por oferta, siempre que pertenezcan al ámbito de Lógica. Al inicio del ciclo lectivo, la persona docente debe comunicar por escrito al estudiantado. \\ \hline
\textbf{MA-0832}\newline Tópicos de Topología (4) & Profundizar en un tema de la topología, definido para cada oferta del curso, mediante el estudio guiado de resultados, problemas y, cuando corresponda, aplicaciones. & Los contenidos no se fijan de manera permanente en el plan de estudios: son abiertos por oferta, siempre que pertenezcan al ámbito de Topología. Al inicio del ciclo lectivo, la persona docente debe comunicar por escrito al estudiantado. \\ \hline
\textbf{MA-0833}\newline Tópicos de Estadística (4) & Profundizar en un tema de la estadística, definido para cada oferta del curso, mediante el estudio guiado de resultados, problemas y, cuando corresponda, aplicaciones. & Los contenidos no se fijan de manera permanente en el plan de estudios: son abiertos por oferta, siempre que pertenezcan al ámbito de Estadística. Al inicio del ciclo lectivo, la persona docente debe comunicar por escrito al estudiantado. \\ \hline
\textbf{MA-0928}\newline Procesos Estocásticos (4) & Manejar los conceptos y técnicas básicas de la teoría de procesos estocásticos en tiempo continuo, con especial atención al movimiento browniano y sus aplicaciones. & Procesos gaussianos; Movimiento browniano; Filtraciones, martingalas y semimartingalas continuas; Variación cuadrática e integración estocástica; Fórmula de Itô y consecuencias; Procesos de Markov; Movimiento browniano y ecuaciones en derivadas parciales; Ecuaciones diferenciales estocásticas. \\ \hline
\end{longtable}
\endgroup
